Para validar la usabilidad de Thary, se siguió el plan de validación ya definido. La realización de la retroalimentación fue hecha por el desarrollador del sistema y por una persona encargada de la gestión de inventarios en el dominio de salud, esto con el fin de obtener una perspectiva externa a la del equipo de desarrollo.

Para un usuario registrado como administrador el flujo del sistema se pudo completar sin dificultad e incluyó: inicio de sesión, carga y procesamiento del historial de consumo, consulta del dashboard de demanda pronosticada, del estado del inventario, del presupuesto y del calendario de pedidos y confirmación (y reversión) de un pedido sugerido. 

Una vez completado todo el proceso, la interfaz fue percibida como intuitiva y fácil de usar, no fue necesario que existiera alguna explicación adicional más allá de la que la propia interfaz del sistema proporcionaba a la hora de completar cada una de las tareas.

Como resultado de esta validación, se aportó una mejora de la usabilidad del sistema. El gestor de inventario que probó el sistema sugirió una mejora concreta: incorporar  una plantilla de archivo CSV descargable desde la propia interfaz que le sirviera a los usuarios de guía sobre el formato y columnas esperadas antes de cargar su propio histórico de consumo.  Esta sugerencia se identificó como valiosa porque, si bien Thary ya estaba diseñado para aceptar distintas variantes de nombres de columna, y en la pantalla de carga se incluye un manual con los nombres de las columnas, la incorporación de una plantilla facilita aún más este proceso al mostrarle al usuario el formato esperado antes de que intente cargar su propio archivo.

Esta sugerencia se implementó en el sistema incluyendo el archivo \texttt{plantilla\_thary.csv}, el cual está disponible para descarga desde la interfaz de carga de archivo.

Se reconoce que esta validación tuvo un carácter informal y con un número reducido de participantes (dos personas), no obstante, al incluir a una persona con experiencia real en gestión de inventario en el dominio de salud, la retroalimentación obtenida se considera representativa del tipo de usuario final para el cual Thary fue diseñado, y se identifica como una línea de trabajo futuro la realización de una validación de usabilidad más extensa, con un número mayor de participantes.